The principal comparison starts with FP32 queries in host memory and ends
when top-10 IDs and scores are available in host memory after stream
completion (\cref{fig:3-boundary}). It includes pinned staging, H2D,
query preprocessing, search and D2H. CAGRA int8 query quantisation is inside
this boundary; quantiser training and index construction are outside it.
Resource handles and output buffers are reused. JHQ returns to pageable
host arrays and the baselines to pinned arrays; CAGRA/PQ use Python APIs,
whereas JHQ/RaBitQ use C++. These interface differences remain visible when
interpreting small margins.

\par\medskip\noindent\begin{minipage}{\linewidth}
\centering\captionsetup{hypcap=false}
\resizebox{\linewidth}{!}{% Protocol diagram, not a measured stage-time plot.
\begin{tikzpicture}[
 font=\small\linespread{1}\selectfont,
 box/.style={draw=ThesisNavy!65,rounded corners=2pt,align=center,
 text width=2.55cm,minimum height=1.55cm,inner sep=4pt},
 arr/.style={-{Latex[length=2mm]},draw=ThesisNavy!80,thick}]
\node[box] (host) at (0,0) {Host FP32 queries\\Pinned staging};
\node[box,fill=ThesisNavy!7] (upload) at (3.5,0) {H2D transfer\\Query preprocessing};
\node[box,fill=ThesisNavy!7] (gpu) at (7,0) {GPU search\\Selection / refinement};
\node[box] (result) at (10.5,0) {D2H IDs and scores\\Stream completion};
\draw[arr] (host) -- (upload);
\draw[arr] (upload) -- (gpu);
\draw[arr] (gpu) -- (result);
\draw[ThesisNavy,thick] (-1.4,1.15) -- (-1.4,1.45) -- (11.9,1.45) -- (11.9,1.15);
\node[fill=white,inner sep=3pt] at (5.25,1.45) {Primary host-to-host timed interval};
\node[box,text width=5.7cm,minimum height=1.3cm] (setup) at (1.65,-2.35)
 {Outside search timing\\Load data / train or load model / add / warm-up};
\node[box,text width=5.7cm,minimum height=1.3cm] (cal) at (8.65,-2.35)
 {Separate policy cost\\Calibration; reuse cost model};
\draw[densely dashed,ThesisNavy!60] (-1.4,-1.1) -- (11.9,-1.1);
\end{tikzpicture}
}
\captionof{figure}[Host and device timing boundaries]{Timing boundaries. The host interval includes transfers and
query-specific computation. Device-stage profiles explain parts of that
interval; separate profiling and synchronisation can change execution, so
their sum is not an interchangeable end-to-end latency.}
\label{fig:3-boundary}
\end{minipage}\par\medskip

With $t$ in seconds, round throughput and its summary are
\begin{equation}
Q_r=\frac{m_rb}{\sum_{i=1}^{m_r}t_{ri}},\qquad
Q=\operatorname{median}_r Q_r,
\label{eq:3-throughput}
\end{equation}
E1 uses batch 1,000 and five rounds. JHQ warms up once per variant/round; baselines warm up three times initially. Normally each round has 24 timed calls, fewer for slow baseline points. Warm-up excludes JIT and capture. Latency percentiles describe repeated batch calls, not individual or concurrent requests.

We retain slow valid calls and both round-QPS and call-latency summaries. Their medians need not be reciprocal. Shared indexes and query pools do not provide independent training replications, so sweep cells receive no independence-based significance claim.

Runs use an RTX 5090 (32,607\,MiB), driver 595.71.05, CUDA 13.0 and cuVS 26.08.01. OpenMP uses 16 threads within a 25-CPU quota; E1 records 575\,W. A cooperative lock and snapshots do not fix clocks or host load. JHQ JIT-compiles Release compute\_80 PTX \cite{cuda_platform}; baselines use installed builds. Device profiles are reported separately.

Resident memory includes encoded arrays, models, workspace and context. E1's context-subtracted JHQ totals and CAGRA's index deltas therefore differ in scope (\cref{sec:4-4,sec:7-7}). Later-created search pools and build temporaries prevent an equal-peak-memory comparison. Separate 20\,ms samples estimate build peaks. Cold train/add excludes input loading and preceding host setup.

\paragraph{Supplementary evidence boundary.}
Supplementary checks (\cref{sec:a-supplement}) use separate sessions. Precision changes and cached-model reuse follow \cref{sec:4-2,sec:7-precision}; historical baseline ratios remain cross-session.